\section{问题三模型建立及求解}

% 后续问题模板：先说明真实继承，再说明问题三新增的对象、条件或数学困难。
% 不复制问题一/二的完整推导；若问题三独立，直接说明新的研究对象和数学缺口。

\subsection{扩展模型的建立}
% 建议顺序：继承的 current 关系 → 新增变量/状态 → 新目标或判据 → 新约束来源 → 最终可计算模型。
% 决定性推导、边界和约束不得因“前文已有模型”而省略；未变化部分用准确引用承接。
% 优化类仍应先说明决策对象和目标函数现实含义，并把目标函数放在约束大括号外。

\subsection{模型求解}
% 从问题三最终模型的计算结构或新增困难进入 solver。
% 写清继承/更换求解方法的原因、本题化编码、约束处理、初值/参数/精度/终止条件和输出映射。
% 算法步骤或伪代码只有在真实多阶段、循环、分支或筛选逻辑需要展示时才启用。

\subsection{模型求解结果与分析}
% 给出 current 直接答案及图表，解释关键结果和有效边界。
% 若继承结果须准确回指；若深化证据修正主结论，则同步回收。
%
% \subsection{模型检验}
% 独立检验确有已验收证据才启用，存在 comparator 时改为“模型检验与对照”。
